The use of artificial intelligence (AI) in healthcare is rapidly increasing. Electronic health records (EHRs) represent a valuable source of data for disease prediction and other machine learning (ML) applications. Because many EHR datasets inherently possess a longitudinal structure, they require specific analytical approaches. Consequently, AI models designed to handle sequential medical data are becoming increasingly common and sophisticated.

Despite these advancements, many AI models trained on EHRs suffer from poor reporting practices, and the lack of clear data presentation significantly limits their practical adoption in clinical and research settings~\cite{placeholder_source}. Furthermore, the majority of machine learning models are currently accessible only via command-line tools or notebook environments, such as Jupyter or Google Colab. While these interfaces are suitable for data scientists and software engineers, they present a steep learning curve for healthcare professionals and real-world evidence specialists who could greatly benefit from models trained on EHR data. 

To bridge this gap, there is a clear need for accessible, user-friendly interfaces. The overarching goal of this thesis is to design and implement an intuitive web application that provides clear visualization of machine learning outputs and enables the analysis of longitudinal administrative patient data. Specifically, the application will serve as a catalog for deploying models focused on predicting selected diseases from sequences of medical codes. These codes are commonly used by Czech healthcare institutions to report procedures to health insurance providers, making the tool highly relevant for real-world healthcare analytics.

The thesis is structured as follows: Chapter~2 provides a review of relevant technologies, focusing on methods for representing longitudinal sequences in EHRs and approaches to integrating ML models into web applications. Chapter~3 details the specific objectives of the thesis based on the official assignment. Chapter~4 describes the system architecture, including the user interface, database structure, and methods for processing model outputs. Chapter~5 presents the results, showcasing the implemented web application and its interactive visualizations. Chapter~6 discusses the outcomes, system expandability, and potential limitations, and Chapter~7 concludes the work.
